\phantomsection\addcontentsline{toc}{section}{Abstract}
\section*{Abstract}
This project investigates how quoted prices respond over time to scheduled Federal Reserve communication and whether those responses support a reliable measure of information-diffusion speed. A reproducible framework combines public event data, one-minute provider bid quotes, unrestricted horizon regressions, synthetic validation and conditional adjustment diagnostics. Scheduled FOMC statements form the primary application; minutes releases and cross-series lead--lag relationships provide exploratory extensions.

The sample covers 2018--2025 and contains 63 scheduled statements and 65 minutes releases. Strict complete-window checks retain 336 of 384 planned event--asset pairs across SPX, NSX and EURUSD provider quotes. A rate-surprise factor is estimated from 39 development statements and held fixed across subsequent samples. Because its measurement window ends 20 minutes after each release, within-window response estimates are retrospective associations rather than announcement-time forecasts or structurally identified policy effects.

A source-clock audit identifies seasonal timestamp behaviour inconsistent with the provider's fixed-EST documentation. Independent minute-quote probes support a deterministic calendar correction, and corrected SPX statement-window returns correlate 0.9949 with published USMPD stock-index returns. The earlier run is preserved. As its later-period outcomes had already been inspected, corrected 2024--2025 results are described as post-audit chronological transfer rather than untouched validation. Primary development and pooled statement coverage still fail the declared 80\% threshold.

For accepted pooled SPX statement windows, the 20-minute response slope is $-29.46$ basis points per development surprise standard deviation. The additional response from minute 5 to minute 20 is $+0.44$ basis points, with a 95\% interval from $-6.66$ to $+7.54$. Minutes releases show a more delayed exploratory association. Neither communication type passes every reporting condition for a single empirical adjustment speed, and later-period prediction-loss comparisons show no convincing improvement over a development-mean baseline.

The principal contribution is an auditable measurement and verification framework with retained negative results, explicit data corrections and interactive companions. Synthetic tests distinguish gradual adjustment from reversals, continuing drift and stale quotes, while also exposing finite-sample uncertainty. The evidence supports conditional response-path measurement; it does not establish the exact time at which all information is incorporated into exchange prices or a profitable trading opportunity.

\textbf{Keywords:} information diffusion; FOMC communication; event studies; price discovery; lead--lag analysis; public data; reproducibility.
